\section{Chunking}
\label{sec:model}
This part of the model admits several possible encodings. A natural but
expensive formulation is a \emph{fully enumerated} encoding, where each chunk
explicitly stores the identifiers of all boxes it contains. Since chunks may
have different lengths, this representation requires a two-dimensional
array with padding values for shorter chunks.
\vspace{0.5cm}
\begin{lstlisting}
array[1..n_chunks, 1..max_pkble_boxes] of var 0..n_boxes: chunks;
\end{lstlisting}
\vspace{0.5cm}
Here, the value \(0\) is used as a dummy value to represent an empty slot.
This representation makes chunk membership explicit and easy to inspect.
However, it introduces a large number of decision variables and requires
additional constraints to enforce the structure of a valid solution: each
box must appear exactly once, and the non-empty positions inside each
chunk must be ordered.

A second possibility is a \emph{binary assignment} encoding, where a boolean
variable states whether box \(i\) belongs to chunk \(g\). This representation
can be written as follows:

\vspace{0.5cm}
\begin{lstlisting}
array[1..n_chunks, 1..n_boxes] of var bool: belongs_to_chunk;
\end{lstlisting}
\vspace{0.5cm}
In this case, chunk membership is represented by the truth value of
\texttt{belongs\_to\_chunk[g,i]}. Although this encoding is expressive, it
is even larger than the fully enumerated representation, and it does not
encode contiguity directly: that property must be reconstructed through
additional constraints.

For this reason, the adopted model represents each chunk using its
\emph{start index} and \emph{length}. An equivalent formulation could use
the start and end indices instead. This interval-style representation is
more compact because each chunk is described by only two variables, and
contiguity is implicit: once the start index and the length are known,
the boxes in the chunk are uniquely determined. As a result, the model
avoids explicit membership variables and reduces the number of constraints
needed to describe valid chunks.

\vspace{0.5cm}
\begin{lstlisting}
array[1..n_chunks] of var 1..n_boxes: chunk_start;
array[1..n_chunks] of var 1..max_pkble_boxes: chunk_len;
\end{lstlisting}
\vspace{0.5cm}
The adopted encoding uses two arrays: \texttt{chunk\_start}, which stores
the first box of each chunk, and \texttt{chunk\_len}, which stores the
number of boxes in that chunk.

The following constraints define the position of each chunk in the input
stream:

\begin{lstlisting}
constraint chunk_start[1] = 1;
\end{lstlisting}

The first constraint sets the start of the first chunk to box \(1\). This fixes
the beginning of the chunk sequence and ensures that the partition starts from
the first box.

\begin{lstlisting}
constraint forall(n in 1..n_chunks - 1)(
  chunk_start[n + 1] = chunk_start[n] + chunk_len[n]
);
\end{lstlisting}

The second constraint makes the chunks consecutive. For each chunk \(n\), the
next chunk starts immediately after it, since the start position of chunk
\(n+1\) is obtained by adding the length of chunk \(n\) to
\texttt{chunk\_start[n]}. This prevents gaps and overlaps between adjacent
chunks.

\begin{lstlisting}
constraint chunk_start[n_chunks] + chunk_len[n_chunks] - 1 = n_boxes;
\end{lstlisting}

The final constraint forces the last chunk to end exactly at box
\texttt{n\_boxes}. Together, these constraints ensure that the chunks form a
complete partition of the input stream, where every box is covered exactly once.